\chapter{文献综述}
\begingroup
\justifying
\setlength{\parindent}{0pt}
\setstretch{1.5}
\setlength{\parskip}{0.5\baselineskip}
\titlespacing{\chapter}{0pt}{0pt}{0pt}
\titlespacing{\section}{0pt}{0pt}{0pt}

\section{电子投票的安全需求}

可信电子投票系统至少应满足四项性质：资格控制、选票保密、计票完整性与可审计性。经典协议文献进一步细化为个体可验证性（选民可确认自身选票被计入）、普遍可验证性（任意观察者可检查计票一致性），以及无收据性与抗胁迫性（两者相关但不等价） \cite{benaloh1994,chaum2004,juels2005}。

在实践中，这些要求分布于密码学控制与运行控制两个层面。密码学可证明某些关系成立，但无法单独保证终端安全、密钥治理或在非受控环境下的抗胁迫行为。解读“安全投票”声明时，这一鸿沟至关重要。

\section{中心化电子投票系统}

传统中心化架构通过 Web 服务与数据库事务处理选举生命周期。成熟工具链与较低运维复杂度赋予其现实优势；但事件响应仍位于同一信任边界内部，其局限同样明显：
\begin{itemize}
    \item 一旦访问控制失守，特权内部人员可能篡改记录或日志；
    \item 若未提供完整数据导出，外部无法独立回放状态迁移；
    \item 争议处理最终常回退到机构信任，而非可验证记录。
\end{itemize}

这些不足推动研究从“信任运营者”转向“信任协议”。

\section{基于区块链的投票方案}

区块链基础设施以确定性执行规则维护不可篡改的追加式日志 \cite{nakamoto2008,wood2014}。选举状态迁移与事件轨迹可透明记录，智能合约可直接强制生命周期不变量，从而减少对隐藏后端逻辑的依赖。

然而，透明性具有双刃剑属性。若投票语义以明文出现，机密性将永久丧失。因此，隐私保护型区块链投票系统必须同时实现：流程控制层面的账本透明，以及选票内容层面的密码学隐藏。

\section{用于隐私计票的同态加密}

椭圆曲线离散对数群上的 ElGamal 加密支持对编码消息的加法同态 \cite{elgamal1985,hankerson2004ecc}。本研究在 BabyJubJub 曲线上实例化该方案；该曲线定义于 BN254 标量域，可直接嵌入 Groth16 电路算术并避免额外域转换开销。每张选票被编码为 one-hot 向量；按分量聚合密文即可得到加密候选人总票数，仅对聚合结果解密可避免暴露单张选票。

该方法具有两项实践优势。其一，密文聚合可通过点加法在线性复杂度下随投票规模扩展；其二，解密点恢复为整数计票可在链下独立执行。只要私钥泄露与明文侧信道风险得到控制，隐私性质即可成立。相较门限解密，其保证更弱，但无需复杂仪式基础设施。

\section{电子投票中的零知识证明}

零知识证明（ZKP）是一类密码学协议：证明者在不泄露除“命题为真”之外任何信息的前提下，说服验证者该命题成立 \cite{goldwasser1989}。在电子投票场景中，这一能力是基础性需求：选民需证明其选票格式正确（且仅编码一个有效候选项）而不泄露所选对象；管理员需证明计票解密正确而不公开私钥。由此，公共可验证性与选票隐私可并存。

本论文采用 Groth16（zk-SNARK）方案，原因在于其配对校验成本相对于电路约束数保持固定，且证明大小恒定为三个椭圆曲线点 \cite{groth2016}。总验证 gas 仍包含较小的公开输入线性项，但不随完整电路规模增长。Circom 与 snarkjs 提供了从电路描述到可部署 Solidity 验证器的完整实现链路 \cite{circomdocs,snarkjs}。

\subsection{ZKP 解决的问题}

许多非零知识证明工作流要求向验证者暴露见证（私密数据）。这在电子投票中不可接受：若为证明选票有效而公开候选人索引，则选票保密性被破坏；若为证明计票正确而公开解密密钥，则密钥治理失效。

ZKP 将\emph{声明}与\emph{证据}分离。设 $x$ 为公开陈述，$w$ 为私密见证。定义关系 $\mathcal{R}$，使得 $(x, w) \in \mathcal{R}$ 表示“$w$ 是 $x$ 的有效见证”。ZKP 系统允许知晓 $w$ 的证明者生成短证明 $\pi$，使任意验证者相信 $(x, w) \in \mathcal{R}$，同时 $\pi$ 不泄露关于 $w$ 的信息。

在本项目中，该模式包含两类实例：
\begin{itemize}
    \item \textbf{投票证明（vote proof）}：公开陈述与见证为
    \begin{align*}
    x &= (\mathit{commitment},\; H_{ct},\; PK_x,\; PK_y) \\
    w &= (\mathit{candidateId},\; \mathit{salt},\; r_0, \dots, r_{N-1},\\
       &\quad \mathit{c1X}_0,\; \mathit{c1Y}_0,\; \mathit{c2X}_0,\; \mathit{c2Y}_0,\; \dots,\; \mathit{c1X}_{N-1},\; \mathit{c1Y}_{N-1},\; \mathit{c2X}_{N-1},\; \mathit{c2Y}_{N-1})
    \end{align*}
    关系 $\mathcal{R}$ 要求：(i) $\mathit{commitment} = \mathsf{Poseidon}(\mathit{candidateId}, \mathit{salt})$；(ii) $0 \le \mathit{candidateId} < N$；(iii) 每组密文 $(\mathit{c1X}_j,\mathit{c1Y}_j,\mathit{c2X}_j,\mathit{c2Y}_j)$ 都是 one-hot 比特 $m_j$ 的有效 ElGamal 加密，且当且仅当 $j=\mathit{candidateId}$ 时 $m_j=1$；(iv) $H_{ct}$ 与全部密文坐标的 Poseidon 哈希一致。
    \item \textbf{计票证明（tally proof）}：令 $w = sk$，公开陈述为
    \begin{equation*}
        \begin{split}
        x = (&PK_x,\; PK_y,\; \bar{C1}^{X}_0,\; \bar{C1}^{Y}_0,\; \dots,\; \bar{C1}^{X}_{N-1},\; \bar{C1}^{Y}_{N-1},\\
             &\bar{C2}^{X}_0,\; \bar{C2}^{Y}_0,\; \dots,\; \bar{C2}^{X}_{N-1},\; \bar{C2}^{Y}_{N-1},\\
             &R^{X}_0,\; R^{Y}_0,\; \dots,\; R^{X}_{N-1},\; R^{Y}_{N-1},\; \mathit{totalVotes})
        \end{split}
    \end{equation*}
    其中 $R_j = n_j \cdot G$ 为候选人 $j$ 的结果点，$R^{X}_j, R^{Y}_j$ 为 $R_j$ 的坐标，$n_j$ 为对应整数票数，$\mathit{totalVotes}$ 为所有候选人票数之和（对应第 1 章中的 $V$）。
其中 $PK_x, PK_y$ 为点 $\mathit{PK}$ 的仿射坐标。
\end{itemize}

\subsection{形式化安全性质}

ZKP 系统须满足三项性质 \cite{goldwasser1989}：
\begin{enumerate}
    \item \textbf{完备性（Completeness）}：若 $(x, w) \in \mathcal{R}$，则诚实证明者总能生成被验证者接受的证明 $\pi$。这保证有效选票与正确计票不会被误拒。
    \item \textbf{可靠性（Soundness）}：若命题 $x$ 不存在有效见证 $w$，则任意计算受限攻击者几乎不可能伪造通过验证的证明。该性质阻止伪造选票与错误计票声明被链上接受。
    \item \textbf{零知识性（Zero-Knowledge）}：存在模拟器可在未知 $w$ 的情况下生成与真实证明同分布的证明。该性质形式化地说明 $\pi$ 不泄露见证细节。
\end{enumerate}

在 zk-SNARK 中，可靠性是计算意义上的（依赖密码学困难假设），非交互性通过可信设置生成的结构化参考字符串（SRS）实现。SRS 提供证明与验证所需参数，无需往返交互。

\subsection{算术电路与 R1CS}

有限域 $\mathbb{F}_p$ 上的任意计算都可表示为由加法门和乘法门组成的算术电路。当存在对全部导线（公开输入、私密见证和中间值）的赋值使所有门方程成立时，电路称为\emph{可满足}。该可满足性问题可编码为\emph{秩一约束系统（R1CS）}。

R1CS 由定义在见证向量 $\mathbf{s} \in \mathbb{F}_p^n$ 上的 $m$ 个约束组成（$\mathbf{s}$ 包含公开输入、私密见证和中间导线值，且约定 $s_0 = 1$）：
\begin{equation}
    \langle \mathbf{a}_i,\, \mathbf{s} \rangle \;\cdot\; \langle \mathbf{b}_i,\, \mathbf{s} \rangle \;=\; \langle \mathbf{c}_i,\, \mathbf{s} \rangle, \qquad i = 1, \dots, m
\end{equation}
其中 $\mathbf{a}_i, \mathbf{b}_i, \mathbf{c}_i \in \mathbb{F}_p^n$ 为由电路拓扑导出的稀疏选择向量。每个乘法门对应一个约束；加法与常数乘法可并入选择向量而不额外增加约束数。

在 circom 中，开发者以 DSL 编写约束，编译器输出 \texttt{.r1cs} 文件（约束矩阵 $A, B, C$）以及 WASM 见证生成器。本项目编译结果为：
\begin{itemize}
    \item \texttt{vote\_proof}：14{,}590 个约束，编码承诺验证、候选人范围校验、one-hot 加密合法性与密文哈希绑定；
    \item \texttt{tally\_proof}：16{,}786 个约束，编码密钥归属验证、逐候选人解密正确性与总票数一致性。
\end{itemize}

\subsection{QAP 变换}

Groth16 在 R1CS 之上使用多项式编码形式，即\emph{二次算术程序（QAP）} \cite{gennaro2013}。对每个约束索引 $i$ 与见证索引 $j$，通过拉格朗日插值定义多项式 $u_j, v_j, w_j \in \mathbb{F}_p[X]$，满足 $u_j(i) = a_{i,j}$、$v_j(i) = b_{i,j}$、$w_j(i) = c_{i,j}$。

定义
\begin{equation*}
A(X)=\sum_j s_j\,u_j(X), \qquad B(X)=\sum_j s_j\,v_j(X), \qquad C(X)=\sum_j s_j\,w_j(X).
\end{equation*}
在每个约束点 $X=i$ 处，由插值可得 $A(i)=\langle \mathbf{a}_i,\mathbf{s}\rangle$、$B(i)=\langle \mathbf{b}_i,\mathbf{s}\rangle$、$C(i)=\langle \mathbf{c}_i,\mathbf{s}\rangle$。因此 R1CS 方程推出
\begin{equation*}
A(i)\cdot B(i)-C(i)=0, \qquad i=1,\dots,m.
\end{equation*}
故多项式
\begin{equation}
\begin{split}
    P(X) \;=\; &\Bigl(\sum_j s_j\, u_j(X)\Bigr) \cdot \Bigl(\sum_j s_j\, v_j(X)\Bigr) \\
               &-\; \sum_j s_j\, w_j(X)
\end{split}
\end{equation}
在全部 $i=1,\dots,m$ 上满足 $P(i)=0$。由于 $T(X)=\prod_{i=1}^{m}(X-i)$ 恰具有这些根，故 $T(X)$ 整除 $P(X)$。等价地，存在商多项式 $H(X)$ 使得：
\begin{equation}
    A(X) \cdot B(X) \;-\; C(X) \;=\; H(X) \cdot T(X)
\end{equation}
其中 $A,B,C$ 分别表示左输入、右输入与输出线性组合。证明者任务是在未知秘密点 $\tau$ 明文值的情况下，对各多项式在 $\tau$ 处的取值进行承诺；这一能力由可信设置阶段产出的 SRS 提供。

\subsection{Groth16 协议}

Groth16 通过椭圆曲线双线性配对实现 QAP 论证。设 $\mathbb{G}_1, \mathbb{G}_2$ 为素数阶 $p$ 的群，生成元分别为 $G_1, G_2$，配对映射为 $e: \mathbb{G}_1 \times \mathbb{G}_2 \rightarrow \mathbb{G}_T$。协议包含三阶段。

\paragraph{Setup.}
可信方采样秘密参数 $(\tau, \alpha, \beta, \gamma, \delta) \overset{\$}{\leftarrow} \mathbb{F}_p^5$，并计算两类结构化参考字符串：
\begin{itemize}
    \item \textbf{证明密钥} $\mathit{pk}_{\mathit{Groth}}$：包含所需幂次的 $\{[\tau^i]_1\}$、$\{[u_j(\tau)]_1\}$、$\{[v_j(\tau)]_2\}$ 及其盲化组合（记号上有别于 ElGamal 公钥 $\mathit{PK}$ 及其坐标 $PK_x,PK_y$）。对本规模电路而言，点数量可达数万级。
    \item \textbf{验证密钥} $\mathit{vk}$：包含 $[\alpha]_1$、$[\beta]_2$、$[\gamma]_2$、$[\delta]_2$，以及公开信号索引对应的 $\{[(\beta u_j + \alpha v_j + w_j)(\tau)/\gamma]_1\}$。
\end{itemize}

随后必须不可恢复地销毁 $(\tau, \alpha, \beta, \gamma, \delta)$，即可信设置中的“有毒废料（toxic waste）”。若任一值泄露，攻击者可在未知有效见证时伪造可通过验证的证明。因此实践中采用 \emph{Powers of Tau} 多方仪式：多个独立参与者各自注入随机性并证明其已销毁本地贡献。只要至少一方诚实，最终产物即保持安全。

\paragraph{Prove.}
给定完整见证 $\mathbf{s}$ 与证明密钥，证明者采样随机盲化标量 $r, s' \overset{\$}{\leftarrow} \mathbb{F}_p$，并计算三元证明 $\pi = (\pi_A, \pi_B, \pi_C)$：
\begin{align}
    \pi_A &= \bigl[\alpha + A(\tau) + r\delta\bigr]_1 \\
    \pi_B &= \bigl[\beta + B(\tau) + s'\delta\bigr]_2 \\
    \pi_C &= \left[\frac{C_{\mathit{priv}}(\tau) + H(\tau)T(\tau)}{\delta}\right]_1 + s' \cdot \pi_A \\
          &\quad + r \cdot \bigl[\beta + B(\tau) + s'\delta\bigr]_1 - rs' \cdot [\delta]_1 \notag
\end{align}
其中 $[\cdot]_k$ 表示椭圆曲线点 $(\cdot) \cdot G_k$，$C_{\mathit{priv}}$ 表示 $C$ 多项式中与私密见证相关部分。盲化因子 $r$ 与 $s'$ 保证 Groth16 的完美零知识性质：模拟器无需见证即可生成与真实证明同分布的样本。

\paragraph{Verify.}
验证者检查一条双线性配对方程，将证明元素与验证密钥元素组合：
\begin{equation}
    e(\pi_A,\; \pi_B) \;=\; e\bigl([\alpha]_1,\; [\beta]_2\bigr) \;\cdot\; e\!\left(\sum_{i=0}^{l} a_i \cdot \mathit{vk}_i,\; [\gamma]_2\right) \;\cdot\; e(\pi_C,\; [\delta]_2)
\end{equation}
其中 $l$ 为显式公开输入数（不含常数项），$\mathit{vk}_i = [(\beta u_i + \alpha v_i + w_i)(\tau)/\gamma]_1$ 为预计算验证密钥项，$a_0 = 1$ 为常数项，$a_1,\dots,a_l$ 为公开输入值。该单一方程利用配对代数性质编码了 $A(\tau) \cdot B(\tau) = C(\tau) + H(\tau) \cdot T(\tau)$ 的正确性，而无需在 $\tau$ 处重算多项式。方程成立当且仅当证明被接受。

\subsection{链上验证成本}

BN254 曲线（亦称 BN128）在 EVM 中通过地址 \texttt{0x08} 的预编译合约原生支持配对运算。根据 EIP-1108 \cite{eip1108}，gas 成本模型为：
\begin{equation}
    \mathit{gas}_{\mathit{pairing}} = 45{,}000 + 34{,}000 \times k
\end{equation}
其中 $k$ 为配对对数。Groth16 验证方程需要 $k = 4$ 对（对应方程四项），因此仅配对校验即消耗 $45{,}000 + 34{,}000 \times 4 = 181{,}000$ gas。关键在于该配对成本\emph{与电路约束数量无关}：14{,}590 约束电路与 100{,}000 约束电路的配对 gas 相同。公开输入线性组合仅引入与公开输入数量相关的小幅开销，而非与总约束相关。这一性质使 Groth16 特别适合需要链上验证的复杂 ZKP 计算。

\subsection{实现流水线}

本实现采用四阶段流程：
\begin{enumerate}
    \item \textbf{电路定义}：使用 circom 编写约束，调用 BabyJubJub 算术、Poseidon 哈希与条件逻辑模板。编译器输出 \texttt{.r1cs}（约束矩阵 $A,B,C$）与 WASM 见证生成器。
    \item \textbf{可信设置}：Powers of Tau 产物（\texttt{pot16\_final.ptau}）提供支持最多 $2^{16}$ 约束的通用 SRS。随后使用 snarkjs 执行电路特定的第二阶段设置，生成证明密钥（\texttt{.zkey} 文件；本实现中每电路约 7--10 MB）。
    \item \textbf{证明生成}：给定公开输入与私密见证，WASM 生成器构造完整见证向量 $\mathbf{s}$；snarkjs 再基于证明密钥计算 $(\pi_A, \pi_B, \pi_C)$。本项目 Node.js 测试显示，投票证明生成约 3--8 秒；浏览器端 snarkjs WASM 通常相近或因设备硬件略慢。
    \item \textbf{验证器导出}：snarkjs 生成内嵌验证密钥的 Solidity 合约，并以 EVM 预编译调用直接实现配对校验方程，避免运行时解析开销。
\end{enumerate}

该流水线的重要信任属性在于\emph{电路代码公开可审计}：任何人可审阅 \texttt{vote\_proof.circom} 与 \texttt{tally\_proof.circom}、重新编译，并在部署时对照公开的 \texttt{.zkey} 产物检查 R1CS 一致性。结合可验证的 Powers of Tau 仪式记录，系统信任基础落在密码学困难假设与设置仪式上，而非电路作者个人诚信。

\subsection{ZKP 在本文中的角色}

本系统部署三类证明实例，组织于两份电路文件：
\begin{enumerate}
    \item \textbf{ZKP1（承诺正确性）}：证明 $\mathit{commitment} = \mathsf{Poseidon}(\mathit{candidateId}, \mathit{salt})$，且不泄露 $\mathit{candidateId}$ 与 $\mathit{salt}$。
    \item \textbf{ZKP2（one-hot 合法性）}：证明提交的 ElGamal 密文确实加密了与承诺候选项一致的 one-hot 向量 $\mathbf{m}$，且不泄露 $\mathbf{m}$ 与分量随机数。由于 ZKP1 与 ZKP2 共享见证变量（均依赖 $\mathit{candidateId}$），两者被编译为同一 \texttt{vote\_proof} 电路，单个证明即可同时成立。
    \item \textbf{ZKP3（计票解密正确性）}：证明对每个候选人 $j$，公开结果点 $\mathbf{results}[j] \cdot G$ 等于在私钥 $sk$ 下对聚合密文解密得到的点，且 $sk$ 与公钥 $PK = sk \cdot G$ 一致。整数票数通过离散对数恢复在链下计算并单独提交；该电路证明的是曲线点表示的正确性，而非直接证明整数值。该电路必须独立存在，因为其输入依赖投票结束后的聚合密文状态。
\end{enumerate}

由 snarkjs 自动生成的链上验证器（\texttt{VoteVerifier.sol}、\texttt{TallyVerifier.sol}）分别以纯 Solidity 实现配对校验方程。主合约 \texttt{Voting} 在 \texttt{castVote} 与 \texttt{updateTallyResults} 中将其作为强制前置条件，从而在状态更新前执行电路关系约束。第 5 章将讨论仍需加固的绑定缺口。

\section{承诺、哈希与 Merkle 审计性}

在本设计中，承诺值是对隐藏选票数据的绑定型回执。选民对候选人 ID 与盐值求哈希，仅公开摘要。后续若任一私密值被修改，摘要会变化，因此初始声明无法被静默改写。此外，由于盐值 $s$ 均匀随机且不公开，承诺同时满足隐藏性：仅凭摘要无法推断 $\mathit{candidateId}$。

承诺步骤采用 Poseidon，原因在于其在零知识电路中约束开销低 \cite{poseidon2019}。这一选择具有直接工程意义：同一哈希需在证明生成中重复验证，若哈希电路成本过高将显著拖慢证明速度。

随后，系统通过 Merkle 树将“单个承诺可验证”扩展为“整组选票可验证” \cite{merkle1988}。每个叶子存储一个承诺，根值绑定整份有序列表。包含性证明长度随叶子数按对数增长，即使选票规模扩大，验证成本仍可控。

由此形成两类独立检查：选民可验证自身承诺是否在公开根下被包含，且验证步骤不额外暴露身份信息；外部审计者可按确定性排序规则重建根值，并与公布根及链上根进行一致性比对。

\section{本研究定位}

本论文实现并评估了完整工程路径，将浏览器侧证明生成、合约级验证与选后审计工具连接为一个可部署工作流。

贡献可归纳为两点：
\begin{itemize}
    \item 提供可运行原型，证明上述机制可在同一架构内协同工作；
    \item 对安全目标达成状态进行显式评估，区分“已完整满足”“部分满足”与“尚未解决”。
\end{itemize}

表 \ref{tab:approach-comparison} 总结了与本论文最相关的架构权衡。

\begin{table}[htbp]
\centering
\small
\caption{不同投票架构的比较}
\label{tab:approach-comparison}
\begin{tabular}{L{3.0cm} L{2.9cm} L{2.9cm} L{3.9cm}}
\toprule
评价维度 & 中心化电子投票 & 仅区块链投票 & 本研究（区块链 + HE + ZKP + Merkle） \\
\midrule
状态透明性 & 受运营方日志公开范围限制 & 链上动作透明度高 & 控制流程与公开承诺均具高透明性 \\
选票机密性 & 依赖后端策略 & 链上明文时通常较弱 & 通过加密选票与承诺回执显著增强 \\
计票可验证性 & 对运营可信性依赖较重 & 未引入密码证明时仅中等 & 通过投票/计票证明提供高可验证性 \\
选民侧包含性校验 & 罕见或与账户强关联 & 可实现但隐私敏感 & 基于回执的 Merkle 包含性验证，降低验证阶段地址关联 \\
运行复杂度 & 较低 & 中等 & 更高，但带来更强保障属性 \\
\bottomrule
\end{tabular}
\end{table}

\section{本章小结}

本章将论文工作置于既有电子投票研究脉络中，并说明仅依赖区块链透明性不足以满足本问题场景需求。要同时获得强隐私与可信计票，必须在设计中引入并整合额外密码学机制。

文献综述同时揭示了核心权衡：所提架构提高了实现复杂度，但这种复杂度换来了更强的可验证性与完整性保证。

\endgroup

